\documentclass[class=report, 12pt, crop=false, oneside]{standalone}

% --- PREAMBLE UNTUK TIMES NEW ROMAN & FORMATTING ---
\usepackage[utf8]{inputenc}
\usepackage[T1]{fontenc}
\usepackage[indonesian]{babel}
\usepackage{graphicx}
\usepackage{enumitem}
\usepackage{amsmath}
\usepackage{mathptmx} % Times New Roman
\usepackage{titlesec}
\usepackage{setspace}
\usepackage[top=2.5cm, bottom=2.5cm, left=4cm, right=3cm]{geometry}
% Memastikan penomoran sampai 4 tingkat (3.2.1.1)
\setcounter{secnumdepth}{4}
\setcounter{tocdepth}{4}

% Memulai dokumen dari Bab 3
\setcounter{chapter}{2} 

% --- FORMAT JUDUL ---
% Judul Bab Tetap di Tengah (Standar Skripsi)
\titleformat{\chapter}[display]
  {\normalfont\bfseries\filcenter}{\chaptertitlename\ \thechapter}{10pt}{\large}

% Section (3.1) Rata Kiri
\titleformat{\section}
  {\normalfont\large\bfseries}{\thesection}{1em}{}

% Subsection (3.2.1) Rata Kiri
\titleformat{\subsection}
  {\normalfont\normalsize\bfseries}{\thesubsection}{1em}{}

% Subsubsection / Paragraph diformat menjadi 3.2.1.1 Rata Kiri
\titleformat{\subsubsection}
  {\normalfont\normalsize\bfseries}{\thesubsubsection}{1em}{}

\onehalfspacing
\graphicspath{{/home/veinmahzy/skrispi/images/bab_3/}} 

\begin{document}
\fontsize{12pt}{14.4pt}\selectfont

\chapter{METODE PENELITIAN}

%================================================================================
% 3.1 KERANGKA BERPIKIR
%================================================================================
\section{Kerangka Berpikir}

\indent Berikut ini gambar Kerangka Berpikir yang digunakan dalam Pengembangan Sistem Pelaporan Masalah Lingkungan Dengan Integrasi Data Gambar dan Teks untuk mendukung \textit{Smart City}:

\begin{figure}[htbp]
    \centering
    \includegraphics[width=0.75\textwidth]{Kerangka Bepikir.png}
    \caption{Kerangka Berpikir Penelitian}
    \label{fig:kerangka_berpikir}
\end{figure}

\indent Penelitian ini bertujuan mengembangkan sistem klasifikasi otomatis yang dapat mengidentifikasi dan mengklasifikasikan setiap laporan masyarakat yang masuk sebagai 8 kelas yaitu jalanan rusak, kecelakaan, kerusakan trotoar, vandalisme, sampah, kerusakan halte, kecelakaan lalu lintas, dan pohon tumbang menggunakan teknologi \textit{Artificial Intelligence} (AI). 

%================================================================================
% 3.2 METODE PENELITIAN
%================================================================================
\section{Metode Penelitian}

Penelitian ini bertujuan untuk mengintegrasikan teknologi \textit{Deep Learning} ke dalam platform pelaporan masyarakat guna mewujudkan sistem klasifikasi otomatis yang presisi. Metodologi penelitian dibagi menjadi dua fase krusial: pengembangan model \textit{Artificial Intelligence} (AI) dan perancangan sistem berbasis \textit{web}. 

%--------------------------------------------------------------------------------
% 3.2.1 Pengembangan Sistem Artificial Intelligence (AI)
%--------------------------------------------------------------------------------
\subsection{Pengembangan Sistem \textit{Artificial Intelligence} (AI)}

Pengembangan AI yang dilakukan oleh peneliti dapat dilihat pada Gambar \ref{fig:ai_development}. Pertama-tama kami melakukan pengumpulan, akuisisi, dan pembersihan dataset.

\begin{figure}[htbp]
    \centering
    \includegraphics[width=0.9\textwidth]{Artifical Intelegence Alur Eskperimen.png}
    \caption{Alur Pengembangan Sistem \textit{Artificial Intelligence}}
    \label{fig:ai_development}
\end{figure}

%--------------------------------------------------------------------------------
% 3.2.1.1 Dataset
%--------------------------------------------------------------------------------
\subsubsection{Dataset}

Tahap pengumpulan dataset melibatkan dua metode perolehan data yaitu penggunaan dataset publik dari platform Kaggle yang berfokus pada deteksi kerusakan jalan, dan pengambilan data secara mandiri melalui metode \textit{web scraping} pada portal pengaduan masyarakat crm.jakarta.go.id. Tim peneliti melakukan \textit{web scraping} terhadap situs tersebut untuk mengumpulkan dataset berisi laporan masyarakat sepanjang tahun 2025 yang bersifat publik.

\end{document}